\chapter{人が耳を傾けたくなる話し方}

この講義は短い。だが、その組み立てには並外れた精密さがある。Treasure はまず、問題の賭け金を大きくする。人間の声は、私たちの誰もが奏でている楽器であり、おそらく世界で最も強力な音なのだ、と。ついで彼は、その壮大さを実際的な謎へと絞り込む。もし媒体としての声がそれほど強いのなら、なぜ話しても聞いてもらえないことが、これほどありふれているのか。答えは、意図的な順序で展開される。まず、聞きやすさを損なう習慣の診断。つぎに、積極的な倫理的基盤。つぎに、声の変数からなる道具箱。つぎに、ウォームアップの手順。そして最後に、意識的な音に満ちた世界とはどのようなものかという、より広い構想である。

\section{声と問題}

冒頭の主張は、意図的に大きい。声は戦争を始めることもできるし、``愛している'' と告げることもできる。Treasure がまず感じさせたいのは、話すことは装飾ではない、ということだ。それは帰結をもつ。そうして十分に拡大したのちにはじめて、彼はこの講演を支配する問いを置く。なぜ人はこれほどしばしば耳を傾け損なうのか。そして、話すことはいかにして、世界を変えるだけの力を持ちうるのか。

簡潔な編集上の略記を用いれば、この逆説は次のように要約できる。
\begin{equation}
\text{強力な媒体} \not\Rightarrow \text{有効なコミュニケーション}.
\end{equation}
声には莫大な潜在力がある。だが、潜在力だけでは、受け取られることは保証されない。

Treasure の次の一手が重要である。彼は、よりよく聞こえるための小技からは始めない。まず、そもそも何が聞くことを妨げるのかを問う。したがってこの講義は、処方に先立って診断から進んでゆく。

\subsection{質問 \& 答え}

\paragraph{質問。} なぜ、私たちが話しても人は耳を傾けないのか。

\paragraph{答え。} Treasure の答えには二つの層がある。第一に、技法の問題に入る以前に、私たちを聞きにくい存在にしてしまう反復的な習慣がある。第二に、意図が善くても、届け方は重要である。同じ言葉でも、どう話されるかによって、まったく異なる強さ、権威、意味を帯びて届く。ゆえに講義は、習慣から基盤へ、基盤から声の力学へと進むのである。

\section{話し方の七つの大罪}

Treasure の最初の答えは、完全に否定的である。彼は ``話し方の七つの大罪'' の一覧を組み立てる。それは完結した道徳体系としてではなく、聞きやすさを蝕む習慣の実際的な分類としてである。その順序自体が講義のリズムの一部をなしているので、明示的に保っておく価値がある。
\begin{equation}
S_7=
\bigl(
\text{陰口},
\text{裁き},
\text{否定性},
\text{不平},
\text{言い訳},
\text{誇張},
\text{独断}
\bigr).
\end{equation}

\begin{remark}
この記法は編集上のものである。講義の秩序立ったレパートリーを紙面に保つためだけのものだ。
\end{remark}

では、講演がそうするように、この一覧を順に追っていこう。

\begin{enumerate}
\item \textbf{陰口。} その場にいない人の悪口を言うことは、ただちに信頼を弱める。話し手が他人の陰口を言うなら、その人は後になって私たちのことも同じように言うだろう、と私たちは自然に推測する。
\item \textbf{裁き。} 同時に裁かれ、欠けた存在と見なされていると感じながら、誰かの話に耳を傾けるのは難しい。裁く話し手は、聞き手を自己防衛へ追い込む。
\item \textbf{否定性。} Treasure はこれを、10月1日はひどい日だという一言で例示する。要点は単なる笑いではない。執拗な否定性は、聞き手の世界を縮こまらせ、言葉を重くする。
\item \textbf{不平。} 否定性の別の形として導入されるが、あまりにありふれているため、それだけで一項目を成している。Treasure はこれを、伝染する惨めさと呼ぶ。光ではなく不幸を広げるのだ。
\item \textbf{言い訳。} ここでの典型は、人に責任を投げつける人間、つまり責任をつねに他人へ回してしまう人である。そうした話し方が逃げているように感じられるのは、責任がどこにも落ち着かないからである。
\item \textbf{誇張。} embroidery によって彼が意味しているのは、誇張である。言葉の較正が失われる。何もかもを awesome と呼んでしまえば、言葉はもはや本当の驚嘆を記録できない。そして誇張は、十分に押し進められれば、嘘になる。
\item \textbf{独断。} 独断とは、事実と意見を取り違えることである。この二つが混同されると、聞き手はもはや思考へ招き入れられているのではなく、真理として扱われた断言に打ち据えられることになる。
\end{enumerate}

ここで重要なのは、累積的な構造である。Treasure は単に「もう少し親切であれ」と言っているのではない。話し手から聞き手への回路が、どのように一つひとつの習慣によって損なわれるかを示しているのだ。一覧の終わりまで来ると、診断は鋭くなる。人が耳を傾けないのは、しばしば、聞くことそのものを難しくしているのが私たち自身だからである。

だからこそ、ここで講義には本当の転回が必要になる。診断のあとには、基盤が求められる。

\section{HAIL: 積極的な基盤}

ここで Treasure は、前節が準備してきた問いを立てる。これを積極的に考える道はあるのか。ある。答えは、\(\mathrm{HAIL}\) という記憶術に圧縮された、もう一つの秩序立った構造である。
\begin{align}
H &\mapsto \text{honesty（正直さ）},\\
A &\mapsto \text{authenticity（真正さ）},\\
I &\mapsto \text{integrity（誠実さ）},\\
L &\mapsto \text{love（愛）}.
\end{align}

この語が重要なのは二つの意味においてである。それは記憶に残りやすい。しかも、挨拶すること、あるいは熱意をこめて迎えることの響きも帯びている。この四つの基盤から発せられた言葉は、そのように受け取られるだろう、という示唆がそこにある。

\subsection{質問 \& 答え}

\paragraph{質問。} 七つの罪に代わる、積極的な基盤とは何か。

\paragraph{答え。} Treasure の答えは、力ある話し方は、真実であり、本物であり、頼るに足り、善意に満ちていなければならない、というものである。正直さは線をまっすぐに保つ。真正さとは、仮面越しに話していないということである。誠実さとは、私たちが自分の言葉そのものであるということだ。愛とは、聞き手の幸福を願うことである。これらは技法のあとに付け足される飾りではない。技法が立つことのできる地面なのである。

このブロックで最も興味深い局所的緊張は、正直さと愛との関係に現れる。Treasure は、それだけで十分であるかのように ``絶対的な正直さ'' を勧めたりはしない。彼はすぐに、不必要な無遠慮さを反例として差し出す。簡潔に言えば、
\begin{equation}
\text{絶対的な正直さ} \not\Rightarrow \text{よい話し方},
\qquad
\text{愛によって和らげられた正直さ} \Rightarrow \text{よい話し方}.
\end{equation}
愛は真実性を打ち消すのではない。その用い方に制約を与えるのである。

ついで彼は、もう一つの、よりやわらかな主張を加える。
\begin{equation}
\text{愛} \;\Rightarrow\; \text{裁きたがる傾向の減少}.
\end{equation}
これは定理ではなく、そのように読まれるべきでもない。質的な両立しにくさを述べているのである。もし私たちが本当に相手の幸福を願っているなら、その人を裁くという姿勢を保つのは難しくなる。

このブロックを Treasure は、決定的な転換句で終える。以上が、私たちが何を語るかである。では、それをどう語るのかを問うことが残っている。

\section{声の道具箱}

ここで講義は、倫理的基盤から、声の力学へと転じる。声は Treasure の言うところでは、ほとんどの人が一度も開いたことのない、驚くべき道具箱である。ノートのためには、この道具をひとつの対象にまとめておくと便利である。
\begin{equation}
V_{\text{tool}}=
\bigl(
\text{声の置き場所},
\text{声色},
\text{プロソディ},
\text{速度},
\text{沈黙},
\text{高さ},
\text{音量}
\bigr).
\end{equation}

要点は、Treasure がこのような公式を書くということではない。彼は書かない。要点は、講義そのものがいまや、話すことを、少数の主要変数をもつ制御可能な系として扱っている、ということだ。

\paragraph{声の置き場所。} Treasure はまず、声がどこに置かれているかを対比する。鼻の高いところ、より中央の喉、そしてさらに低い胸である。彼が提供しているのは、形式的な音響学の講義ではなく、話すことの実践的な現象学である。重みがほしいなら、人は下へ降りる傾向がある、と彼は言う。そこからすぐに、彼の修辞的で経験的な主張が導かれる。
\begin{equation}
\text{低い声域} \;\Rightarrow\; \text{より大きな力と権威として知覚される}.
\end{equation}
だから彼は、私たちは声の低い政治家に投票するのだ、と言うのである。要点は数値測定ではない。知覚される力なのである。

\paragraph{声色。} 声色とは、声がどのように感じられるかである。Treasure は、望ましい声を、豊かで、滑らかで、暖かい、``ホットチョコレートのような'' 声として描く。この比喩のなかに隠れている技術的な点は重要である。声色は訓練可能なのだ。呼吸、姿勢、運動によって、それは変えられる。

\paragraph{プロソディ。} Treasure にとって、プロソディは会話における意味への第一の道筋である。それは、話し言葉が力と方向を得るための、歌うようなメタ言語である。だからこそ、単調な話し方は聞き取りにくい。それは、意味を運ぶ主要な担い手の一つを取り去ってしまうからだ。

\subsection{質問 \& 答え}

\paragraph{質問。} 同じ言葉が、話し方によって違う意味を持ちうるのはなぜか。

\paragraph{答え。} 意味は、語の列そのものの中だけに住んでいるわけではないからである。語尾を持ち上げる反復的な調子の例や、``Where did you leave my keys?'' という一文の例が示すのは、言葉が固定されていても、輪郭が発話行為そのものを変えてしまうということだ。

この考えを、最小限の形式で書いてみよう。いま
\begin{align}
u &= \text{``Where did you leave my keys''},\\
m_1 &= \mathcal{M}(u,\pi_1),\\
m_2 &= \mathcal{M}(u,\pi_2),
\end{align}
とする。ここで \(\pi_1\) と \(\pi_2\) は異なるプロソディないし音高の輪郭であり、\(\mathcal{M}\) は聞き手に知覚される意味である。すると Treasure の要点は、ただ
\begin{equation}
m_1 \neq m_2.
\end{equation}
ということに尽きる。

\begin{remark}
これもまた編集上の記法である。実演の論理、すなわち同じ言葉でも、異なる輪郭によって異なる意味になるという点を捉えるためのものだ。
\end{remark}

だからこそ、語尾を上げる調子の反復は問題なのである。すべての文が問いであるかのように上昇して終わると、平叙文は平叙の力を失う。一つの輪郭だけでは、意図された意味の全域を担えない。

\paragraph{速度と沈黙。} ついで Treasure は、対比によって速度を実演する。非常に速く話せば興奮を生み出せるし、ぐっと遅くすれば強調を生み出せる。その尺度の果てにあるのが沈黙である。沈黙は、話し続けられなかった失敗ではない。伝達のための道具なのだ。
\begin{equation}
\text{速い速度} \Rightarrow \text{興奮},
\qquad
\text{遅い速度} \Rightarrow \text{強調},
\qquad
\text{沈黙} \Rightarrow \text{力ある間}.
\end{equation}
だからこそ、話は ``um'' や ``ah'' で埋め尽くされる必要はない。沈黙には重みを持たせることができる。

\paragraph{高さ。} 音の高さは、しばしば速度と組み合わさって覚醒や昂揚を示すが、Treasure は、それが単独でも働きうることを示す。同じ文を二通りに発したとき、その違いは旋律的輪郭の違いから生まれる。言葉は固定されている。変えられているのは意図である。

\paragraph{音量。} 最後の道具は音量である。大きな声は、指令力や興奮を生みうる。小さな声は、部屋全体を身を乗り出させることができる。Treasure の警告は、絶え間ない拡声に向けられている。彼はこれを ``sodcasting'' と呼ぶ。つまり、他人に対して不注意に音を押しつけることだ。

この節全体は、実際に技術的な転換を遂げている。私たちはもはや、ただ ``うまく話しなさい'' と言われているのではない。意図的に変化させることのできる変数の集合を示されており、その一つひとつが、聞き手が発話をどう受け取るかを変えるのである。

\section{楽器を準備する}

次の転換は、動機づけの転換である。Treasure は、話すことが本当に重大になる場面を考えよ、と促す。公の講演、提案、昇給の願い、結婚式のスピーチ。そうした場面では、自分自身に対して準備する義務があるのだ、と彼は言う。ここで比喩はさらに鋭くなる。声は単なる道具箱ではなく、エンジンでもあり、どんなエンジンも、暖まらずにうまくは動かない。

\subsection{質問 \& 答え}

\paragraph{質問。} なぜ、話す前に声をウォームアップするのか。

\paragraph{答え。} 大事な発話には、楽器への配慮がふさわしいからである。もし声が、コミュニケーションという行為を運ぶエンジンなのだとすれば、冷えたまま始動するべきではない。Treasure はこれを比喩の水準にとどめない。彼はそれを、具体的な本番前の手順へと変える。

簡潔な参照のために、次のように書いてもよい。
\begin{equation}
W_6=
\bigl(
\text{呼吸とため息},
\text{唇の打音},
\text{リップトリル},
\text{舌の } la\text{ 練習},
\text{巻き舌の } r,
\text{サイレン } (we \to or)
\bigr),
\end{equation}
だが、本当の内容は順序立った手順のなかにある。
\begin{enumerate}
\item 両腕を上げ、深く息を吸い、ため息とともに吐き出す。
\item ``bop'' という音を繰り返して、唇を温める。
\item 子どもがぶるぶると鳴らすような音、すなわちリップトリルへ進む。
\item 誇張した ``la, la, la'' の発音で、舌を目覚めさせる。
\item \(r\) を巻いて発音する。Treasure はこれを印象的に ``舌にとってのシャンパン'' と呼ぶ。
\item 高い ``we'' から低い ``or'' へ移るサイレンで締めくくる。
\end{enumerate}

この手順が重要なのは、それが講義における最初の、完全に実践的な見返りだからである。Treasure は、悪徳から基盤へ、基盤から変数へ、そしていまや手順へと進んできた。聴衆は、何が重要かを告げられるだけでなく、それを実際に稽古するよう求められているのである。

そして彼は、このブロック全体に短い命令的コーダを付け加える。次に話すときには、前もってそれをやりなさい。この一行は、説明が習慣へ変わる印だから、ノートの中でも見える形で残しておくべきである。

\section{結びの構想: 意識的に話すこと、意識的に聞くこと、意識的な音}

講義の終わりで Treasure は、この議論全体を文脈の中に置きたいのだと言う。この動きは彼らしい。コンパクトなレパートリーと局所的実演の連鎖のあとで、彼は視野を広げ、こうした実践がどのような世界を築くのかを問う。

現在の状態についての彼の診断は、簡潔には次のように書ける。
\begin{equation}
\text{現在の世界}
=
\bigl(
\text{貧しい話し方},
\text{貧しい聴き方},
\text{騒音},
\text{悪い音響}
\bigr).
\end{equation}
これは個人の欠陥であるだけではない。環境の欠陥でもある。人はまずく話し、聞き手は耳を傾けず、音響環境そのものが理解に逆らっている。

彼が私たちに思い描かせたい反事実的な世界は、その反対である。
\begin{equation}
\text{望ましい世界}
=
\bigl(
\text{力ある話し方},
\text{意識的な聴取},
\text{意識的な音の創造},
\text{目的にかなった音響}
\bigr).
\end{equation}

したがって講義は、私的な自己改善の教訓で終わるのではない。市民的な音響学で終わるのである。もし私たちが意識的に音を作り、意識的に音を受け取り、意識的に音のために環境を設計したなら、何が起こるだろうか。Treasure の答えは、そのような世界は美しく響き、理解が標準となるだろう、というものだ。

この最後の拡大は、講演のリズムにとって決定的である。人がなぜ耳を傾けないのかという問いとして始まったものは、私たちが話す習慣と聞く習慣によって、どのような人間世界を築いているのかという問いとして終わる。

\section{まとめ}

Treasure の講義は、それぞれが前のものによって動機づけられた、入れ子状の構造の連なりからできている。私たちはまず、強力な声がしばしば失敗するという逆説から始める。ついで、その失敗を七つの罪によって診断する。HAIL によって、道徳的な地盤を立て直す。声の道具箱を開き、届け方は偶然ではなく変数によって左右されることを見る。その楽器を六つのウォームアップによって準備する。そして最後に、探究全体を意識的な音の世界へと広げて終える。

したがってこの章は、話し方のこつを詰め込んだ袋としてではなく、コミュニケーションについての小さくも構造だった理論として読むべきである。話すことが重要なのは、それが力を持つからだ。それは、同定可能な理由によって失敗する。それは、明示的な基盤の上に立て直すことができる。その力は、制御可能な変数に依存している。そして、私たちがそれらの変数を意図的に使うことを学ぶとき、得られるのは修辞的な有効性だけではない。理解という、より大きな可能性そのものなのである。